\documentclass[
	12pt
] {article}

\usepackage[
	a4paper,
	left=3.5cm,
	right=1.5cm,
	top=3cm,
	bottom=3.2cm,
	headheight=1cm,
	footskip=1.6cm
] {geometry} %for setting margin

\usepackage{fancyhdr} %for header and footers
\usepackage{parskip} %for leaving some space after paragraphs


% Math
\usepackage{mathtools} %helps to number equations
\usepackage{siunitx} % Units of measurement


% For graphs and advanced diagrams
\usepackage{tikz} %the drawing library
\usetikzlibrary{positioning} %some auxilary to the drawing library
\usetikzlibrary{math} %helps to define variables in the tikz
\usetikzlibrary{patterns} %helps to define variables in the tikz
\usetikzlibrary{patterns.meta} %helps to define variables in the tikz
\usetikzlibrary{hobby}

% Basic graphics
\usepackage{float} %to force the figure to be in a place among the text
\usepackage{graphicx} %for include graphics
\usepackage{caption} %better referencing to labels
\usepackage{subcaption}  %helps with multiple pictures in a figure


% Language and Font package: COMMENT THESE FOR ENGLISH ONLY
\usepackage[T2A]{fontenc} %[RUS] loads cyrillic characters
\usepackage[russian]{babel} %[RUS] changes date heading 'abstract' to Russian
\usepackage[utf8]{inputenc} %[RUS] to force utf-8 encoding
\usepackage{hyperref}

%-----------------------------------------------------------------------
\tikzset{
    every path/.style={
        line width=0.8pt % New default thickness
    }
}
\usetikzlibrary{patterns.meta}


\definecolor{cw}{rgb}{0.4, 0.4, 0.4}
\definecolor{cdw}{rgb}{0.4, 0.4, 0.4}

\definecolor{cnw}{rgb}{0.8, 0.8, 0.8}
\definecolor{cdnw}{rgb}{0.8, 0.8, 0.8}

\definecolor{cbg}{rgb}{0.95, 0.9, 0.5}

\begin{document}

	\newpage \begin{figure}[H]
		\centering
		\begin{subfigure}[b]{\textwidth}
			\centering
			\begin{tikzpicture}
	\tikzmath{
		\widthpic = 8;
		\small = \widthpic / 400; %the filling rectangle in the nodes should be a bit submerged
		\l = 4; %distance between the nodes
	}

	%CENTRAL NODE PARAMETERS
	\tikzmath{
		\r = 0.8;
	}

	% draw central node
	\draw[fill=cnw] (0, 0) circle (\r);

	\begin{scope}[rotate=45]
		%characterstics of 2nd node and properties
		\tikzmath{
			\d = 0.6; %thickness of the capillary
		}
		\tikzmath{
			\pb = sqrt(\r * \r - \d * \d / 4); % p border
			\pr = \pb - \small; %p length
		}
		\begin{scope}[shift={(\pr, 0)}]
			%default one
			\fill[cnw] (0, \d/2) rectangle (\l, -\d/2);
			%draw the main outer borders
			\draw[] (0, \d/2) -- (\l, \d/2);
			\draw[] (0, -\d/2) -- (\l, -\d/2);
		\end{scope}
	\end{scope}


	\begin{scope}[rotate=135]
		%characterstics of 2nd node and properties
		\tikzmath{
			\d = 0.4; %thickness of the capillary
		}
		\tikzmath{
			\pb = sqrt(\r * \r - \d * \d / 4); % p border
			\pr = \pb - \small; %p length
		}
		\begin{scope}[shift={(\pr, 0)}]
			%default one
			\fill[cnw] (0, \d/2) rectangle (\l, -\d/2);
			%draw the main outer borders
			\draw[] (0, \d/2) -- (\l, \d/2);
			\draw[] (0, -\d/2) -- (\l, -\d/2);
		\end{scope}
	\end{scope}

	%tube
	\begin{scope}[rotate=-45]
		%characterstics of 2nd node and properties
		\tikzmath{
			\d = 0.8; %thickness of the capillary
		}
		\tikzmath{
			\pb = sqrt(\r * \r - \d * \d / 4); % p border
			\pr = \pb - \small; %p length
		}
		\begin{scope}[shift={(\pr, 0)}]
			%default one
			\fill[cnw] (0, \d/2) rectangle (\l, -\d/2);
			%draw the main outer borders
			\draw[] (0, \d/2) -- (\l, \d/2);
			\draw[] (0, -\d/2) -- (\l, -\d/2);
		\end{scope}
	\end{scope}




	\begin{scope}[rotate=-135]
		%characterstics of 2nd node and properties
		\tikzmath{
			\d = 0.8; %thickness of the capillary
		}
		\tikzmath{
			\pb = sqrt(\r * \r - \d * \d / 4); % p border
			\pr = \pb - \small; %p length
		}
		\begin{scope}[shift={(\pr, 0)}]
			%default one
			\fill[cw] (0, \d/2) rectangle (\l, -\d/2);

			%draw a slob of rectangle

			\begin{scope}
				\tikzmath{
					\pos1 = 0;
					\pos2 = 0.2;
				}
				\fill[cnw] (\pos1 * \l, \d/2) rectangle (\pos2 * \l, -\d/2);
			\end{scope}

			%drawing a meniscus
			\begin{scope}
				%SET MENISCUS POSITION
				\tikzmath{
					\meniscuspos = 0.2; %in 0-1 coordinates
					\meniscusrot = 0; %0=) 180=(
				}
				%drawing meniscus
				\begin{scope}[rotate=\meniscusrot, shift={(\meniscuspos*\l - \d/4,0)}]
					\fill[cw] (0, \d / 2) rectangle (\d/2, -\d / 2);
					\fill[cnw] (0, -\d/2) arc(-90:90:\d /2) -- cycle;
					\draw[] (0, -\d/2) arc(-90:90:\d /2);
				\end{scope}
			\end{scope}

			%draw the main outer borders
			\draw[] (0, \d/2) -- (\l, \d/2);
			\draw[] (0, -\d/2) -- (\l, -\d/2);

		\end{scope}
	\end{scope}


\end{tikzpicture}

		\end{subfigure}
		\caption{Sample picture generated from tikz}
	\end{figure}


	\newpage \begin{figure}[H]
		\centering
		\begin{subfigure}[b]{\textwidth}
			\centering
			\begin{tikzpicture}
\begin{scope}[rotate=-90]
	\tikzmath{
		\widthpic = 8;
		\small = \widthpic / 400; %the filling rectangle in the nodes should be a bit submerged
		\l = 4; %distance between the nodes
	}

	%CENTRAL NODE PARAMETERS
	\tikzmath{
		\r = 0.8;
	}

	% draw central node
	\draw[fill=cnw] (0, 0) circle (\r);
	\node[] at (0, 0) {$\nu_0$};
	\begin{scope}[rotate=45]
		%characterstics of 2nd node and properties
		\tikzmath{
			\d = 0.6; %thickness of the capillary
		}
		\tikzmath{
			\pb = sqrt(\r * \r - \d * \d / 4); % p border
			\pr = \pb - \small; %p length
		}
		\begin{scope}[shift={(\pr, 0)}]
			\draw[dashed] (\l + \r/2.5, 0) circle (\r/1.5);
			%default one
			\fill[cnw] (0, \d/2) rectangle (\l, -\d/2);


			%draw the main outer borders
			\draw[] (0, \d/2) -- (\l, \d/2);
			\draw[] (0, -\d/2) -- (\l, -\d/2);
			\draw[->, thick] (0.1*\l, 0) -- (0.5*\l, 0);
			\node[] at (0.25*\l, -0.6) {$Q_1$};
			\node[] at (1.05*\l, 0) {$\nu_1$};
		\end{scope}
	\end{scope}


	\begin{scope}[rotate=135]
		%characterstics of 2nd node and properties
		\tikzmath{
			\d = 0.4; %thickness of the capillary
		}
		\tikzmath{
			\pb = sqrt(\r * \r - \d * \d / 4); % p border
			\pr = \pb - \small; %p length
		}
		\begin{scope}[shift={(\pr, 0)}]
			\draw[dashed] (\l + \r/4, 0) circle (\r/2);
			%default one
			\fill[cnw] (0, \d/2) rectangle (\l, -\d/2);




			%draw the main outer borders
			\draw[] (0, \d/2) -- (\l, \d/2);
			\draw[] (0, -\d/2) -- (\l, -\d/2);
			\draw[->, thick] (0.1*\l, 0) -- (0.2*\l, 0);
			\node[] at (0.25*\l, -0.6) {$Q_2$};
			\node[] at (1.05*\l, 0) {$\nu_2$};
			\draw[] (0.8 * \l, 0) -- (0.8 * \l, 2*\d) node[left] {$2$};
		\end{scope}
	\end{scope}

	%tube
	\begin{scope}[rotate=-45]
		%characterstics of 2nd node and properties
		\tikzmath{
			\d = 0.8; %thickness of the capillary
		}
		\tikzmath{
			\pb = sqrt(\r * \r - \d * \d / 4); % p border
			\pr = \pb - \small; %p length
		}
		\begin{scope}[shift={(\pr, 0)}]
			\draw[dashed] (\l + \r/2.5, 0) circle (\r/1.5);
			%default one
			\fill[cnw] (0, \d/2) rectangle (\l, -\d/2);

			%~ %draw a slob of rectangle
			%~ \begin{scope}
				%~ \tikzmath{
					%~ \pos1 = 0;
					%~ \pos2 = 0.4;
				%~ }
				%~ \fill[cdnw] (\pos1 * \l, \d/2) rectangle (\pos2 * \l, -\d/2);
			%~ \end{scope}


			%draw the main outer borders
			\draw[] (0, \d/2) -- (\l, \d/2);
			\draw[] (0, -\d/2) -- (\l, -\d/2);
			\draw[<-, thick] (0.1*\l, 0) -- (0.4*\l, 0);
			\node[] at (0.25*\l, -0.8) {$Q_3$};
			\node[] at (1.05*\l, 0) {$\nu_3$};
		\end{scope}
	\end{scope}




	\begin{scope}[rotate=-135]
		%characterstics of 2nd node and properties
		\tikzmath{
			\d = 0.8; %thickness of the capillary
		}
		\tikzmath{
			\pb = sqrt(\r * \r - \d * \d / 4); % p border
			\pr = \pb - \small; %p length
		}
		\begin{scope}[shift={(\pr, 0)}]
			\draw[dashed] (\l + \r/2.5, 0) circle (\r/1.5);
			%default one
			\fill[cw] (0, \d/2) rectangle (\l, -\d/2);

			%draw a slob of rectangle
			\begin{scope}
				\tikzmath{
					\pos1 = 0;
					\pos2 = 0.2;
				}
				\fill[cnw] (\pos1 * \l, \d/2) rectangle (\pos2 * \l, -\d/2);
			\end{scope}

			%drawing a meniscus
			\begin{scope}
				%SET MENISCUS POSITION
				\tikzmath{
					\meniscuspos = 0.2; %in 0-1 coordinates
					\meniscusrot = 0; %0=) 180=(
				}
				%drawing meniscus
				\begin{scope}[rotate=\meniscusrot, shift={(\meniscuspos*\l - \d/4,0)}]
					\fill[cw] (0, \d / 2) rectangle (\d/2, -\d / 2);
					\fill[cnw] (0, -\d/2) arc(-90:90:\d /2) -- cycle;
					\draw[] (0, -\d/2) arc(-90:90:\d /2);
				\end{scope}
				\draw[] (\meniscuspos * \l + \d/4, \d/4) -- (\meniscuspos * \l + \d/4, + 2*\d) node[left] {$3$};
			\end{scope}


			%draw the main outer borders
			\draw[] (0, \d/2) -- (\l, \d/2);
			\draw[] (0, -\d/2) -- (\l, -\d/2);
			\draw[<-, thick] (0.1*\l, 0) -- (0.4*\l, 0);
			\node[] at (0.25*\l, -0.8) {$Q_4$};
			\node[] at (1.05*\l, 0) {$\nu_4$};
			\draw[] (0.8 * \l, 0) -- (0.8 * \l, - 2*\d) node[right] {$1$};
		\end{scope}
	\end{scope}
\end{scope}
\end{tikzpicture}

		\end{subfigure}
		\caption{Sample picture generated from tikz}
	\end{figure}

	\newpage \begin{figure}[H]
		\centering
		\begin{subfigure}[b]{\textwidth}
			\centering
			\input{fig-sample-3}
		\end{subfigure}
		\caption{Sample picture generated from tikz}
	\end{figure}

	\newpage \begin{figure}[H]
		\centering
		\begin{subfigure}[b]{\textwidth}
			\centering
			\begin{tikzpicture}
	\tikzmath{
		\widthpic = 8;
		\small = \widthpic / 400; %the filling rectangle in the nodes should be a bit submerged
		\l = 4; %distance between the nodes
	}

	%CENTRAL NODE PARAMETERS
	\tikzmath{
		\r = 0.8;
	}

	% draw central node
	\draw[fill=cw] (0, 0) circle (\r);

	\begin{scope}[rotate=45]
		%characterstics of 2nd node and properties
		\tikzmath{
			\d = 0.6; %thickness of the capillary
		}
		\tikzmath{
			\pb = sqrt(\r * \r - \d * \d / 4); % p border
			\pr = \pb - \small; %p length
		}
		\begin{scope}[shift={(\pr, 0)}]
			%default one
			\fill[cnw] (0, \d/2) rectangle (\l, -\d/2);



			%~ \begin{scope}
				%~ \tikzmath{
					%~ \pos1 = 0;
					%~ \pos2 = 0.2;
				%~ }
				%~ \fill[cnw] (\pos1 * \l, \d/2) rectangle (\pos2 * \l, -\d/2);
			%~ \end{scope}

			%draw a slob of rectangle
			\begin{scope}
				\tikzmath{
					\pos1 = 0;
					\pos2 = 0.2;
				}
				\fill[cdnw] (\pos1 * \l, \d/2) rectangle (\pos2 * \l, -\d/2);
			\end{scope}


			%drawing a meniscus
			\begin{scope}
				%SET MENISCUS POSITION
				\tikzmath{
					\meniscuspos = 0.03; %in 0-1 coordinates
					\meniscusrot = 180; %0=) 180=(
				}

				%drawing meniscus
				\begin{scope}[shift={(\meniscuspos*\l,0)}]
					\begin{scope}[rotate=\meniscusrot]
						\begin{scope}[shift={(-\d/4,0)}]
							\fill[cw] (0, \d / 2) rectangle (\d/2 + \small, -\d / 2);
							\fill[cdnw] (0, -\d/2) arc(-90:90:\d /2) -- cycle;
							\draw[] (0, -\d/2) arc(-90:90:\d /2);
						\end{scope}
					\end{scope}
				\end{scope}
			\end{scope}
			%draw the main outer borders
			\draw[] (0, \d/2) -- (\l, \d/2);
			\draw[] (0, -\d/2) -- (\l, -\d/2);

			%\draw[->, thick] (0.1*\l, 0) -- (0.5*\l, 0);
			%\node[] at (0.25*\l, -0.6) {$Q_1$};
		\end{scope}
	\end{scope}


	\begin{scope}[rotate=135]
		%characterstics of 2nd node and properties
		\tikzmath{
			\d = 0.4; %thickness of the capillary
		}
		\tikzmath{
			\pb = sqrt(\r * \r - \d * \d / 4); % p border
			\pr = \pb - \small; %p length
		}
		\begin{scope}[shift={(\pr, 0)}]
			%default one
			\fill[cnw] (0, \d/2) rectangle (\l, -\d/2);

			%draw a slob of rectangle
			\begin{scope}
				\tikzmath{
					\pos1 = 0;
					\pos2 = 0.6;
				}
				\fill[cdnw] (\pos1 * \l, \d/2) rectangle (\pos2 * \l, -\d/2);
			\end{scope}

			%drawing a meniscus
			\begin{scope}
				%SET MENISCUS POSITION
				\tikzmath{
					\meniscuspos = 0.03; %in 0-1 coordinates
					\meniscusrot = 180; %0=) 180=(
				}
				%drawing meniscus
				\begin{scope}[shift={(\meniscuspos*\l,0)}]
					\begin{scope}[rotate=\meniscusrot]
						\begin{scope}[shift={(-\d/4,0)}]
							\fill[cw] (0, \d / 2) rectangle (\d/2 + \small, -\d / 2);
							\fill[cdnw] (0, -\d/2) arc(-90:90:\d /2) -- cycle;
							\draw[] (0, -\d/2) arc(-90:90:\d /2);
						\end{scope}
					\end{scope}
				\end{scope}
			\end{scope}

			%draw the main outer borders
			\draw[] (0, \d/2) -- (\l, \d/2);
			\draw[] (0, -\d/2) -- (\l, -\d/2);
			%\draw[->, thick] (0.1*\l, 0) -- (0.2*\l, 0);
			%\node[] at (0.25*\l, -0.6) {$Q_2$};
		\end{scope}
	\end{scope}

	%tube
	\begin{scope}[rotate=-45]
		%characterstics of 2nd node and properties
		\tikzmath{
			\d = 0.8; %thickness of the capillary
		}
		\tikzmath{
			\pb = sqrt(\r * \r - \d * \d / 4); % p border
			\pr = \pb - \small; %p length
		}
		\begin{scope}[shift={(\pr, 0)}]
			%default one
			\fill[cnw] (0, \d/2) rectangle (\l, -\d/2);

			%drawing a meniscus
			\begin{scope}
				%SET MENISCUS POSITION
				\tikzmath{
					\meniscuspos = 0.03; %in 0-1 coordinates
					\meniscusrot = 180; %0=) 180=(
				}
				%drawing meniscus
				\begin{scope}[shift={(\meniscuspos*\l,0)}]
					\begin{scope}[rotate=\meniscusrot]
						\begin{scope}[shift={(-\d/4,0)}]
							\fill[cw] (0, \d / 2) rectangle (\d/2 + \small, -\d / 2);
							\fill[cnw] (0, -\d/2) arc(-90:90:\d /2) -- cycle;
							\draw[] (0, -\d/2) arc(-90:90:\d /2);
						\end{scope}
					\end{scope}
				\end{scope}
			\end{scope}


			%draw the main outer borders
			\draw[] (0, \d/2) -- (\l, \d/2);
			\draw[] (0, -\d/2) -- (\l, -\d/2);
			%\draw[<-, thick] (0.1*\l, 0) -- (0.4*\l, 0);
			%\node[] at (0.25*\l, -0.6) {$Q_3$};
		\end{scope}
	\end{scope}




	\begin{scope}[rotate=-135]
		%characterstics of 2nd node and properties
		\tikzmath{
			\d = 0.8; %thickness of the capillary
		}
		\tikzmath{
			\pb = sqrt(\r * \r - \d * \d / 4); % p border
			\pr = \pb - \small; %p length
		}
		\begin{scope}[shift={(\pr, 0)}]
			%default one
			\fill[cw] (0, \d/2) rectangle (\l, -\d/2);

			%draw the main outer borders
			\draw[] (0, \d/2) -- (\l, \d/2);
			\draw[] (0, -\d/2) -- (\l, -\d/2);
			%\draw[<-, thick] (0.1*\l, 0) -- (0.4*\l, 0);
			%\node[] at (0.25*\l, -0.8) {$Q_4$};
		\end{scope}
	\end{scope}

\end{tikzpicture}

		\end{subfigure}
		\caption{Sample picture generated from tikz}
	\end{figure}


	\newpage \begin{figure}[H]
		\centering
		\begin{subfigure}[b]{\textwidth}
			\centering
			\input{fig-sample-5}
		\end{subfigure}
		\caption{Sample picture generated from tikzdd}
	\end{figure}

	\newpage \begin{figure}[H]
		\centering
		\begin{subfigure}[b]{\textwidth}
			\centering
			\input{fig-sample-6}
		\end{subfigure}
		\caption{Sample picture generated from tikzdd}
	\end{figure}


	\newpage \begin{figure}[H]
		\centering
		\begin{subfigure}[b]{\textwidth}
			\centering
			\input{fig-sample-7}
		\end{subfigure}
		\caption{Sample picture generated from tikzdd}
	\end{figure}

	\newpage \begin{figure}[H]
		\centering
		\begin{subfigure}[b]{\textwidth}
			\centering
			\input{fig-sample-8}
		\end{subfigure}
		\caption{Sample picture generated from tikzdd}
	\end{figure}


	\newpage \begin{figure}[H]
		\centering
		\begin{subfigure}[b]{\textwidth}
			\centering
			\input{fig-sample-9}
		\end{subfigure}
		\caption{Sample picture generated from tikzdd}
	\end{figure}


	\newpage \begin{figure}[H]
		\centering
		\begin{subfigure}[b]{\textwidth}
			\centering
			\input{fig-sample-10}
		\end{subfigure}
		\caption{Sample picture generated from tikzdd}
	\end{figure}

	\newpage \begin{figure}[H]
		\centering
		\begin{subfigure}[b]{\textwidth}
			\centering
			\begin{tikzpicture}

	\fill[cnw] (0,0) rectangle (5, 5);
	\fill[cw, rounded corners=10pt] (0, 0) -- (2, 1) -- (1.5, 1.5) -- (3.5, 3.5) -- (0, 5) -- (0, 0);
	\fill[cnw, rounded corners=1pt] (0.8, 0.8) -- (1.1, 1) -- (1.3, 1.2) -- (1.2, 1.2) -- (0.7, 0.9) -- cycle;
	\fill[cnw, rounded corners=1pt] (0.8, 2.8) -- (1.1, 3) -- (1.2, 3.5) -- (1.0, 3.6) -- (0.6, 3) -- cycle;
	\fill[pattern={north east lines}] (2,0) rectangle (2.5,5);
	\draw (2, -0.5) -- (2, 5.5);
	\draw (2.5, -0.5) -- (2.5, 5.5);
	\draw[<->] (2, 5.3) -- (2.5, 5.3);
	\node at (2.25, 5.8) {$\Delta x$};
	\node at (-1, 0) {$P_{in}(t)$};
	\node at (6, 0) {$P_{out}$};
	\node at (-2, 2.5) {$Q=const$};

	\draw[->, draw=cw, line width=4] (-1, 3.5) -- (0, 3.5);
	\draw[->, draw=cw, line width=4] (-1, 1.5) -- (0, 1.5);
	\draw[->, draw=cnw, line width=4] (5, 1.5) -- (6, 1.5);
	\draw[->, draw=cnw, line width=4] (5, 3.5) -- (6, 3.5);

	\draw[] (0,0) -- (5, 0);
	\draw[] (0,5) -- (5, 5);

	\draw[-, thick] (0.5,4) -- (-0.5,5) node[left] {1};
	\draw[-, thick] (4.5,4) -- (5.5,5) node[right] {2};
\end{tikzpicture}

		\end{subfigure}
		\caption{Sample picture generated from tikzdd}
	\end{figure}
	\[
		Q_1 + Q_2 = Q_3 + Q_4
	\]

	\[
		q = -\frac{k}{\mu} \nabla P
	\]
	\[
		k = k(S), \quad P_c =  P_c(S)
	\]
	\[
		S = S_{wetting}
	\]

	\[
		q = Q / A
	\]

	\begin{equation}
        P_c = P_c(S, \xi), \quad k = k(S, \xi)
    \end{equation}
    \begin{equation} \label{eq:kondaurov-diff-equ}
        \frac{\partial \xi}{\partial t} = \Omega(S, \xi)
    \end{equation}
    \begin{equation} \label{eq:positive-dissipation-of-cap-energy}
        \frac{\partial \xi}{\partial t} = - C\frac{\partial \psi}{\partial \xi}
    \end{equation}
     \begin{equation} \label{eq:kondaurov-helmholtz}
        \psi = \frac{1}{2} \beta \xi^2 + \frac{1}{2} MS_{nw}^2 - \alpha \xi S_{nw} + \gamma S_{nw}
    \end{equation}
    \begin{equation}\label{eq:series}
        \frac{\beta}{\alpha}\xi = (1 - S) + \tau\frac{\partial S}{\partial t} - \tau^2\frac{\partial^2 S}{\partial t^2} + \tau^3\frac{\partial^3 S}{\partial t^3} - \cdots
    \end{equation}
    \begin{equation}
        \eta = S + \tau \frac{\partial S}{\partial t} = 2S + \frac{\beta}{\alpha} \xi - 1
    \end{equation}
    \begin{equation}
        k_w(S, \xi) = k_w^e(2S + \frac{\beta}{\alpha} \xi - 1)
    \end{equation}
    \begin{equation}
        k_{nw}(S, \xi) = k_{nw}^e(2 - 2S - \frac{\beta}{\alpha} \xi)
    \end{equation}

	\newpage
    $S$ -- насыщенность смачивающей жидкости \\
    $k$ -- проницаемость \\
    $P_c$ -- капиллярное давление \\
    $\xi$ -- параметр капиллярной неравновесности \\
    $\psi$ -- свободная энергия Гельмгольца \\
    $e$ -- равновесные фазовые проницаемости \\
    $C$, $\alpha$, $\beta$, $\gamma$, $M$ -- положительные константы \\

\end{document}
